% =====================================================================
% ejercicio04.tex
% Responsabilidad única: evaluar la entropía cruzada, deducir su
% gradiente, ejecutar un paso de descenso de gradiente y analizar la
% penalización L2 y el parámetro C. Cifras verificadas con
% verificacion/teoria.py.
% =====================================================================
\subsection{Ejercicio 4 --- Regresión logística: pérdida, gradiente y regularización}

Se tienen tres observaciones $(x_i,y_i)$: $(-1,0)$, $(0,0)$ y $(1,1)$,
con $p_i=\sigma(\beta_0+\beta_1x_i)$ y la entropía cruzada media
\[
    J = -\frac{1}{3}\sum_{i=1}^{3}\bigl[y_i\log p_i + (1-y_i)\log(1-p_i)\bigr],
\]
partiendo de $\beta_0=\beta_1=0$ y usando logaritmos naturales.

\begin{enumerate}

\item \textbf{Probabilidades y pérdida inicial.}
Con $\beta=0$ todos los puntajes son $z_i=0$, de forma que
$p_1=p_2=p_3=\sigma(0)=0.5$, y cada término de la suma vale $\log0.5$
(si $y_i=1$ aporta $\log p_i$ y si $y_i=0$ aporta $\log(1-p_i)$, ambos
iguales a $\log0.5$), por lo tanto
\[
    J(0,0) = -\frac13\cdot3\log 0.5 = \log 2 \approx 0.6931 ,
\]
que es la pérdida de un modelo que no distingue entre clases y asigna
probabilidad $\tfrac12$ a todo.

\item \textbf{Gradiente de $J$.}
Para una observación, sea
$\ell_i=-[y_i\log p_i+(1-y_i)\log(1-p_i)]$ con $p_i=\sigma(z_i)$ y
$z_i=\beta_0+\beta_1x_i$. Por la regla de la cadena y
$\sigma'(z)=\sigma(z)(1-\sigma(z))$,
\[
    \frac{\partial\ell_i}{\partial z_i}
    = -\Bigl[\frac{y_i}{p_i} - \frac{1-y_i}{1-p_i}\Bigr]p_i(1-p_i)
    = -\bigl[y_i(1-p_i) - (1-y_i)p_i\bigr]
    = p_i - y_i ,
\]
donde los factores $p_i(1-p_i)$ se cancelan limpiamente con los
denominadores (esa cancelación es precisamente la razón por la que la
entropía cruzada y la sigmoide combinan tan bien, como se resume en la
figura~\ref{fig:ej4-grafo}). Como
$\partial z_i/\partial\beta_0=1$ y $\partial z_i/\partial\beta_1=x_i$,
y $J$ es el promedio de los $\ell_i$,
\[
    \frac{\partial J}{\partial \beta_0} = \frac13\sum_i (p_i-y_i),
    \qquad
    \frac{\partial J}{\partial \beta_1} = \frac13\sum_i (p_i-y_i)\,x_i .
\]
En el punto inicial, $p-y=(0.5,\,0.5,\,-0.5)$, de modo que
\[
    \frac{\partial J}{\partial \beta_0} = \frac{0.5}{3} = \frac16 \approx 0.1667,
    \qquad
    \frac{\partial J}{\partial \beta_1} = \frac{-0.5+0-0.5}{3} = -\frac13 \approx -0.3333 .
\]

\begin{figure}[H]
\centering
\small
\begin{tikzpicture}[
    >={Stealth[length=2.2mm]},
    nodo/.style={draw, rounded corners=2pt, minimum height=8mm,
                 inner sep=4pt, fill=white},
    dato/.style={text=gray!70!black},
    atras/.style={->, thick, classred, dashed},
]
    % --- fila superior: cálculo hacia adelante
    \node[nodo, draw=classblue, fill=classblue!8] (b) at (0,0)
        {$\boldsymbol\beta=(\beta_0,\beta_1)$};
    \node[nodo] (z) at (4.4,0) {$z_i=\beta_0+\beta_1x_i$};
    \node[nodo] (p) at (9.0,0) {$p_i=\sigma(z_i)$};
    \node[nodo] (l) at (13.3,0) {$\ell_i(p_i,y_i)$};
    \node[dato] (x) at (4.4,1.25) {$x_i$};
    \node[dato] (y) at (13.3,1.25) {$y_i$};
    \draw[->, thick] (b) -- (z);
    \draw[->, thick] (z) -- (p);
    \draw[->, thick] (p) -- (l);
    \draw[->, gray!70!black] (x) -- (z);
    \draw[->, gray!70!black] (y) -- (l);
    % --- fila inferior: derivadas acumuladas hacia atrás
    \node[classred] (gb) at (0,-2.3) {$(p_i-y_i)\,(1,\;x_i)$};
    \node[draw=classgreen, very thick, fill=classgreen!12, rounded corners=2pt,
          inner sep=4pt] (gz) at (4.4,-2.3)
        {$\dfrac{\partial\ell_i}{\partial z_i}=p_i-y_i$};
    \node[classred] (gp) at (9.0,-2.3)
        {$\dfrac{\partial\ell_i}{\partial p_i}=\dfrac{p_i-y_i}{p_i(1-p_i)}$};
    \draw[atras] (l.south) |- (gp.east)
        node[pos=0.3, left, font=\footnotesize]
        {$-\Bigl[\dfrac{y_i}{p_i}-\dfrac{1-y_i}{1-p_i}\Bigr]$};
    \draw[atras] (gp) -- (gz)
        node[midway, above, font=\footnotesize] {$\times\,\sigma'(z_i)$}
        node[midway, below, font=\footnotesize] {$=p_i(1-p_i)$};
    \draw[atras] (gz) -- (gb)
        node[midway, above, font=\footnotesize] {$\times\,(1,\;x_i)$};
    \draw[gray!50, densely dotted] (b) -- (gb);
    \draw[gray!50, densely dotted] (z) -- (gz);
    \draw[gray!50, densely dotted] (p) -- (gp);
\end{tikzpicture}
\caption{Grafo computacional de la pérdida de una observación. Arriba,
el cálculo hacia adelante desde los parámetros hasta $\ell_i$; abajo
(rojo), la regla de la cadena aplicada de derecha a izquierda, donde
cada arista multiplica por la derivada local (en la última,
$\partial z_i/\partial\beta_0=1$ y $\partial z_i/\partial\beta_1=x_i$). Al pasar de $p_i$ a $z_i$
el factor $\sigma'(z_i)=p_i(1-p_i)$ cancela exactamente el denominador
de $\partial\ell_i/\partial p_i$, de forma que queda el residuo
$p_i-y_i$ (verde), y al promediar $(p_i-y_i)(1,x_i)$ sobre las tres
observaciones se obtiene el gradiente $\nabla J$ del literal~b.}
\label{fig:ej4-grafo}
\end{figure}

\item \textbf{Un paso de descenso de gradiente.}
Con $\eta=0.5$ y actualización simultánea,
\[
    \beta_0^{\text{nuevo}} = 0 - 0.5\cdot\tfrac16 = -\tfrac{1}{12} \approx -0.0833,
    \qquad
    \beta_1^{\text{nuevo}} = 0 - 0.5\cdot\bigl(-\tfrac13\bigr) = \tfrac16 \approx 0.1667 ,
\]
así que los nuevos puntajes son $z=(-0.25,\,-0.0833,\,0.0833)$ y

\begin{table}[H]
\centering
\begin{tabular}{crrrrr}
\toprule
$i$ & $x_i$ & $y_i$ & $z_i$ & $p_i$ & término de pérdida \\
\midrule
1 & $-1$ & 0 & $-0.2500$ & 0.4378 & $-\log(1-0.4378) = 0.5759$ \\
2 & $0$  & 0 & $-0.0833$ & 0.4792 & $-\log(1-0.4792) = 0.6523$ \\
3 & $1$  & 1 & $0.0833$  & 0.5208 & $-\log(0.5208) = 0.6523$ \\
\bottomrule
\end{tabular}
\caption{Probabilidades y pérdidas individuales tras un paso de descenso de gradiente.}
\end{table}

\[
    J^{\text{nuevo}} = \frac{0.5759+0.6523+0.6523}{3} \approx 0.6269 < 0.6931 ,
\]
por lo cual \textbf{$J$ sí disminuyó} (en $0.0663$), lo esperable al
avanzar en la dirección opuesta al gradiente con un paso moderado: el
intercepto bajó (dos de las tres observaciones son de clase $0$) y la
pendiente pasó a ser positiva (la única observación positiva está en
$x=1$), de forma que las probabilidades se mueven en la dirección
correcta para las tres observaciones.

\item \textbf{Penalización $\lambda\beta_1^2/2$.}
Como el término añadido no depende de $\beta_0$ y
$\frac{\partial}{\partial\beta_1}\bigl(\lambda\beta_1^2/2\bigr)=\lambda\beta_1$,
\[
    \frac{\partial J_\lambda}{\partial\beta_0} = \frac13\sum_i (p_i-y_i),
    \qquad
    \frac{\partial J_\lambda}{\partial\beta_1} = \frac13\sum_i (p_i-y_i)\,x_i + \lambda\beta_1 .
\]
Evaluando en $(\beta_0,\beta_1)=(-\tfrac1{12},\tfrac16)$ con $\lambda=1$
y $p\approx(0.4378,\,0.4792,\,0.5208)$:
\begin{align*}
    \frac{\partial J_\lambda}{\partial\beta_0}
    &= \frac{0.4378+0.4792-0.4792}{3} \approx 0.1459,\\
    \frac{\partial J_\lambda}{\partial\beta_1}
    &= \underbrace{\frac{-0.4378+0-0.4792}{3}}_{\partial J/\partial\beta_1\,\approx\,-0.3057}
    + 1\cdot0.1667 \approx -0.1390 ,
\end{align*}
de modo que la penalización reduce a menos de la mitad el empuje para
seguir aumentando $\beta_1$, actuando como una fuerza que lo atrae hacia
$0$. Sin embargo, en el primer paso no altera nada, porque al partir de
$\beta_1=0$ el término extra $\lambda\beta_1$ vale exactamente $0$, y el
gradiente de $J_\lambda$ coincide con el de $J$, así que la
regularización solo empieza a actuar una vez que $\beta_1$ se aleja de
cero.

\item \textbf{$C$, regularización y escala de entrada.}
En \texttt{LogisticRegression} de scikit-learn, $C$ es el
\textbf{inverso de la intensidad de regularización}: valores pequeños de
$C$ penalizan más los coeficientes (modelos más simples, coeficientes
más cercanos a $0$), y valores grandes de $C$ los penalizan menos,
acercándose al ajuste sin regularizar. Cualitativamente, $C$ juega un
papel inverso al de $\lambda$, aunque no conviene suponer que
$C=1/\lambda$ reproduce exactamente $J_\lambda$, porque, con penalización
L2, scikit-learn minimiza $\tfrac12\lVert w\rVert^2+C\sum_{i=1}^{n}\ell_i$,
es decir, escala la pérdida como una \emph{suma} sobre las observaciones
multiplicada por $C$ (y no como un promedio), de forma que al dividir
entre $Cn$ se obtiene $\tfrac1n\sum_i\ell_i+\tfrac{1}{2Cn}\lVert w\rVert^2$,
y la correspondencia exacta con $J_\lambda$ es $C=1/(n\lambda)$ (aquí
$C=1/3$ y no $C=1$), siempre que el intercepto no se penalice. Este
último punto depende del \emph{solver}: \texttt{lbfgs},
\texttt{newton-cg}, \texttt{newton-cholesky}, \texttt{sag} y
\texttt{saga} dejan el intercepto fuera de $w$, mientras que
\texttt{liblinear} lo trata como el coeficiente de una columna
sintética constante (de valor \texttt{intercept\_scaling}) y por lo
tanto también lo penaliza, de forma que con ese solver ni siquiera
$C=1/(n\lambda)$ reproduce $J_\lambda$ (con los tres datos del
ejercicio y $\lambda=1$, el mínimo de $J_\lambda$ es
$(\beta_0,\beta_1)\approx(-0.7025,\,0.2906)$, que coincide con el de
\texttt{LogisticRegression(C=1/3)} con \texttt{lbfgs}, mientras que
$C=1$ da $\beta_1\approx0.7016$ y \texttt{liblinear} con $C=1/3$ da
$\beta_0\approx-0.134$). Por otra parte, la penalización
$\beta_j^2$ actúa sobre la magnitud numérica de cada coeficiente, y esa
magnitud depende de las unidades: si una variable se mide en milímetros
en lugar de metros, su coeficiente sin penalizar se divide entre $1000$,
de forma que una misma penalización la afecta de manera muy distinta.
Así, a diferencia de mínimos cuadrados o logística sin regularizar
(cuyas predicciones no cambian al reescalar las entradas), un ajuste
regularizado \textbf{sí cambia} con la escala, y por eso se estandarizan
los predictores antes de regularizar, para que todos los coeficientes
compitan en igualdad de condiciones (como se hace en el ejercicio 8 con
\texttt{StandardScaler} dentro del \texttt{Pipeline}).

\end{enumerate}
